%% Stage report (IEEEtran, English) — completed work up to 2026-09-12
%% Compile: pdflatex -interaction=nonstopmode stage_report_ieee_en.tex  (twice)
%% Companion: stage_report_ieee_zh.tex (independent, xelatex + xeCJK)
\documentclass[conference]{IEEEtran}

\usepackage{amsmath}
\usepackage{amssymb}
\usepackage{booktabs}
\usepackage{graphicx}
\usepackage{float}
\usepackage[T1]{fontenc}
\usepackage{lmodern}
\usepackage{microtype}
\usepackage{url}

\begin{document}

\title{Reproduction of DeReFusion and Diagnostics of Operator\\ Specialization for Financial Time Series}

\author{\IEEEauthorblockN{Dongxu Jiang}
\IEEEauthorblockA{Email: jiangdongxu04@gmail.com}}

\maketitle

\begin{abstract}
This report consolidates the work completed up to 2026-09-12 on the reproduction of
DeReFusion and on a sequence of controlled diagnostics that ask a single question: does the
structural state of the input determine which operator is more useful? We first reproduce the
core comparison of the original paper --- a DLinear base branch, an LSTM--Transformer residual
branch and parameter-free additive fusion, all under RevIN --- and then stratify forecast errors
by a causal relative-volatility measure. The marginal value of the nonlinear branch turns out to
be \emph{asset-dependent} rather than universal: significantly positive and seed-robust on the
S\&P~500, directionally positive but statistically inconclusive for ETHUSD, and significantly
\emph{reversed} for BTCUSD. A capacity-sensitivity experiment (hidden width $23\!\to\!128$, all
else frozen) shows that the earlier ``linear dominates everywhere'' reading was partly an artefact
of a narrow nonlinear bottleneck: at adequate capacity ETHUSD prefers the nonlinear operator in
$36/36$ observations while BTCUSD never does. A capacity-matched proxy routing experiment (twelve
pre-registered structural states, linear versus MLP, identical protocol) finds no robust
state-level operator specialization. Finally, an exploratory sweep over ten assets yields
interaction effects that span zero (six negative, four positive, five supported) and two
leave-one-out-stable associations (realized volatility and $|\mathrm{ACF}_1|$), reported as a
\emph{candidate} asset-level structural regularity only. We state explicitly what the evidence
excludes, what it supports, and what remains unresolved.
\end{abstract}

\begin{IEEEkeywords}
reproduction, time series forecasting, operator specialization, structural states,
adaptive fusion, controlled comparison.
\end{IEEEkeywords}

\section{Scope and Research Questions}
The work follows an evidence-first protocol: reproduce the reference result, then ask one
question at a time with every grouping rule, threshold and operator definition fixed \emph{before}
outcomes are inspected. Three questions structure the report: \textbf{Q-A} does the published
ordering of fusion strategies reproduce on independently fetched data; \textbf{Q-B} is the
marginal contribution of the nonlinear branch conditioned on the volatility regime of the input
window; and \textbf{Q-C} does the \emph{structural state} of the input determine which operator is
more useful, i.e.\ is there an experimental basis for routing? All experiments run on CPU under a
frozen 70/10/20 chronological split and report paired bootstrap intervals (4000 resamples) with
paired win rates. No hyper-parameter was tuned after observing test outcomes.

\section{Reproduction of the Reference Comparison}
The dataset of the original repository is not redistributable, so ten daily OHLC series
(2016--2025) were re-fetched; row counts match the paper (2513 for indices and equities, 2602 for
FX, 3653 for BTCUSD, 2975 for ETHUSD). Table~\ref{tab:repro} reports the S\&P~500 comparison at
$T\in\{1,24\}$ (seed 2021).

\begin{table}[!t]
\caption{Reproduction of the core comparison (GSPC, seed 2021, normalized scale).}
\label{tab:repro}
\centering
\small
\begin{tabular}{@{}lcccc@{}}
\toprule
Model & $T{=}1$ & $T{=}24$ & $R^2$ & Params \\
\midrule
DeReFusion (additive)   & 0.01647 & \textbf{0.06230} & \textbf{0.8691} & 28{,}436 \\
revin-DLinear (linear)  & \textbf{0.01577} & 0.07007 & 0.8528 & 4{,}664 \\
gatev1 (volatility)     & ---     & 0.06767 & 0.8578 & 28{,}580 \\
gatev2 (learnable)      & 0.02526 & 0.07191 & 0.8489 & $\sim$28k \\
\bottomrule
\end{tabular}
\end{table}

At $T=24$ the published ordering reproduces: additive fusion beats the linear branch by $11.1\%$
in MSE, while both gate variants are worse than parameter-free addition. At $T=1$ the linear
branch is marginally better, consistent with the behaviour of long-horizon linear forecasters.
Extending the protocol to two further assets at $T=24$ gives $0.21797$ versus $0.22322$ for BTCUSD
($+2.4\%$) and $0.14602$ versus $0.14876$ for ETHUSD ($+1.8\%$); the additive advantage is ordered
GSPC $\gg$ BTCUSD $\approx$ ETHUSD. Figure~\ref{fig:curves} shows the corresponding
per-horizon prediction curves and the run dashboard of the full model.

\begin{figure}[!t]
\centering
\includegraphics[width=0.49\columnwidth]{figs/fig_derefusion_t24_curves.png}
\hfill
\includegraphics[width=0.49\columnwidth]{figs/fig_derefusion_t24_dashboard.png}
\caption{Reproduction of DeReFusion on GSPC ($T{=}24$): forecast curves (left) and the
per-run diagnostic dashboard (right).}
\label{fig:curves}
\end{figure}

\section{Volatility-Regime Stratification}
Each test window is assigned a regime from \emph{causal} volatility measures computed on the input
window only: the standard deviation of log returns and a de-trended version,
$RV^{rel}_t = RV_t / \mathrm{median}(RV_{t-120:t-1})$. Table~\ref{tab:regime} reports the relative
improvement of the nonlinear branch per stratum (relative volatility, the primary measure) and the
interaction effect $\Delta_{high50}-\Delta_{low50}$.

\begin{table}[!t]
\caption{Stratified performance, relative volatility
($\Delta=$ MSE$_{nonlin}-$MSE$_{lin}$; negative favours the nonlinear branch).
$^\dagger$CI excludes zero.}
\label{tab:regime}
\centering
\small
\begin{tabular}{@{}lcccc@{}}
\toprule
Asset & low 20\% & high 50\% & high 20\% & Interaction \\
\midrule
GSPC (2021) & $+6.0\%$ & $+14.2\%^\dagger$ & $+21.6\%^\dagger$ & $-0.01205^\dagger$ \\
GSPC (2022) & $-2.7\%$ & $+23.2\%^\dagger$ & $+29.6\%^\dagger$ & $-0.02248^\dagger$ \\
ETHUSD      & $+16.7\%^\dagger$ & $+6.3\%^\dagger$ & $+14.2\%^\dagger$ & $-0.00943$ \\
BTCUSD      & $+3.5\%$ & $-1.0\%$ & $-3.5\%$ & $+0.01502^\dagger$ \\
\bottomrule
\end{tabular}
\end{table}

Two findings follow. The S\&P~500 effect is \emph{seed-robust}: both seeds give significantly
negative interactions of the same order, and the calm strata are indistinguishable between
branches. BTCUSD reverses the sign significantly, so ``nonlinear advantage grows with volatility''
cannot be stated as a universal criterion. A methodological check explains part of the earlier
confusion: raw volatility is strongly confounded with time (correlation of $RV$ with the sample
index: $+0.44$ for GSPC, $-0.67$ for BTCUSD). The de-trended $RV^{rel}$ removes that trend, and
under it the spurious ``linear wins in the middle quintile'' of GSPC disappears; all reported
figures therefore use $RV^{rel}$.

\section{Adaptive Fusion Does Not Pay}
The natural use of the regime structure would be a gate $\alpha(X)$ mixing the two branches, and
the framework already provides exactly this variant. On the asset where the structure is strongest
and seed-robust (GSPC, $T=24$) the volatility-aware gate obtains MSE $0.06767$, i.e.\ $8.6\%$
\emph{worse} than parameter-free addition ($0.06230$) and only $3.4\%$ better than the plain linear
branch ($0.07007$). Because the nonlinear branch is never worse than the linear one in the calm
strata, the optimal gate weight is essentially constant at one, leaving no switching headroom.
Together with the previously reported negative results for the learnable gate, this motivates
stopping the adaptive-fusion line rather than building a router on top of it.

\section{Capacity Sensitivity}
The proxy experiment of Section~\ref{sec:routing} compares a linear operator with a
capacity-matched MLP ($9{,}240$ versus $9{,}431$ parameters, $+2.1\%$). That matching forces a
384-dimensional input through a 23-dimensional hidden layer, a severe bottleneck. We therefore
re-ran the structural-state evaluation with the hidden width as the \emph{only} variable
($23\to64\to128$; data, features, splits, optimizer, learning rate, batch size, epochs, early
stopping, seeds and evaluation protocol unchanged). Table~\ref{tab:cap} reports the mean
$\Delta$MSE over twelve structural states and three seeds.

\begin{table}[!t]
\caption{Capacity sensitivity. Mean $\Delta$MSE (negative $=$ nonlinear better); in parentheses,
the number of the 36 (state, seed) observations preferring the nonlinear operator.}
\label{tab:cap}
\centering
\small
\begin{tabular}{@{}lcccc@{}}
\toprule
Hidden & Params & GSPC & BTCUSD & ETHUSD \\
\midrule
23  & 9{,}431  & $+0.0598$ (0/36) & $+0.1783$ (0/36) & $+0.0145$ (12/36) \\
64  & 26{,}200 & $+0.0250$ (7/36) & $+0.1888$ (0/36) & $\mathbf{-0.0281}$ (36/36) \\
128 & 52{,}376 & $+0.0061$ (8/36) & $+0.2998$ (0/36) & $\mathbf{-0.0321}$ (36/36) \\
\bottomrule
\end{tabular}
\end{table}

Three pre-registered checks agree. (i) \emph{Sign reversals}: 25 state-level reversals across
widths, 23 of them ETHUSD (all seeds) and 8 GSPC (all at seed 2023 only, i.e.\ seed-unstable);
BTCUSD has none. (ii) \emph{Spread across states}: the range of $\Delta$MSE across states
\emph{shrinks} with capacity for GSPC ($0.106\to0.060$) and ETHUSD ($0.089\to0.040$); the BTCUSD
increase is driven by a single high-error state. (iii) \emph{Per-seed direction consistency} falls
from $31/36$ to $29/36$ to $28/36$. The ETHUSD preference is \emph{asset-wide} once capacity is
sufficient --- all twelve states reverse, not a subset --- so it is asset-level heterogeneity
rather than state-level specialization. We therefore record the earlier uniform ``linear dominates''
statement as partly capacity-driven, and do not treat the capacity sweep as routing evidence.

\section{Structure-Aware Routing: A Minimal Controlled Proxy}
\label{sec:routing}
\subsection{Design}
Twelve structural states were defined \emph{before} any outcome was inspected, using causal
window features: autocorrelation at lags 1/5/10, a trend $t$-statistic, sign persistence, a jump
ratio with a fixed threshold of $c=3\sigma$, the de-trended relative volatility, skewness and
kurtosis. Binary states use train-set medians as thresholds; a composite \emph{structure strength}
state and a $K{=}3$ $k$-means cross-check are also reported. The two operators are a linear map
($9{,}240$ parameters) and an MLP with one 23-unit hidden layer ($9{,}431$ parameters), trained
under an identical protocol (same seed, optimizer, steps, batch, learning rate, early stopping and
preprocessing).

\subsection{Result}
At matched capacity the linear operator is better in \emph{every} one of the twelve states for
GSPC and BTCUSD (GSPC $\Delta$MSE between $+0.047$ and $+0.081$, significant in all three seeds;
BTCUSD between $+0.129$ and $+0.448$). Only ETHUSD shows between-state variation of both signs
(structure-weak states favouring the nonlinear operator by up to $-0.013$, structure-strong states
favouring the linear one by $+0.033$). No state preference survives the capacity sweep.
\textbf{Main finding: no robust operator specialization is established at capacity-matched
scale}; the one cross-asset invariant is that the two operators are indistinguishable in calm
regimes.

\section{Asset-Level Heterogeneity Across Ten Assets}
The frozen protocol ($T=24$, seed 2021) was then run on seven additional assets to map the
\emph{sign} of the interaction effect. GSPC aggregates its two seeds into a single asset, so
$N=10$. Table~\ref{tab:assets} reports the interaction effect and Table~\ref{tab:assoc} the
exploratory association between asset structure features and that effect, including a
leave-one-asset-out (LOO) check.

\begin{table}[!t]
\caption{Asset-level interaction effects (relative volatility).
$^\dagger$CI excludes zero.}
\label{tab:assets}
\centering
\small
\begin{tabular}{@{}lccc@{}}
\toprule
Asset & $\Delta_{interaction}$ & 95\% CI & Supported \\
\midrule
EURUSD & $-0.1371$ & $[-0.164,-0.111]$ & yes$^\dagger$ \\
GSPC   & $-0.0173$ & $[-0.027,-0.009]$ & yes$^\dagger$ \\
USDJPY & $-0.0122$ & $[-0.026,+0.002]$ & no \\
ETHUSD & $-0.0094$ & $[-0.019,+0.000]$ & no \\
DJI    & $-0.0084$ & $[-0.019,+0.003]$ & no \\
TM     & $-0.0081$ & $[-0.028,+0.011]$ & no \\
SOX    & $+0.0044$ & $[-0.008,+0.017]$ & no \\
BTCUSD & $+0.0150$ & $[+0.006,+0.024]$ & yes$^\dagger$ \\
BABA   & $+0.0405$ & $[+0.022,+0.060]$ & yes$^\dagger$ \\
NVO    & $+0.0976$ & $[+0.046,+0.149]$ & yes$^\dagger$ \\
\bottomrule
\end{tabular}
\end{table}

\begin{table}[!t]
\caption{Exploratory Spearman association of asset structure features with
$\Delta_{interaction}$ ($N=10$) and LOO stability.}
\label{tab:assoc}
\centering
\small
\begin{tabular}{@{}lcccc@{}}
\toprule
Feature & $\rho$ & $p$ & LOO range & Stable \\
\midrule
$|\mathrm{ACF}_1|$  & $+0.733$ & $0.016$ & $[+0.633,+0.817]$ & yes \\
Realized volatility & $+0.661$ & $0.038$ & $[+0.533,+0.900]$ & yes \\
Kurtosis            & $+0.345$ & $0.328$ & $[+0.100,+0.533]$ & yes \\
Trend persistence   & $+0.042$ & $0.907$ & $[-0.317,+0.283]$ & no (flips) \\
Jump ratio          & ---      & ---     & --- & constant \\
\bottomrule
\end{tabular}
\end{table}

The sign of the effect spans zero: six assets negative, four positive, five statistically
supported. Two features are LOO-stable: assets with higher realized volatility and stronger
first-order autocorrelation show a \emph{less} favourable (or reversed) nonlinear advantage in
turbulent regimes. At $N{=}10$ both associations reach $p<0.05$ ($|\mathrm{ACF}_1|$
$\rho=+0.733$, $p=0.016$; realized volatility $\rho=+0.661$, $p=0.038$), and the inclusion of
SOX (interaction $+0.0044$, statistically inconclusive) leaves the verdict unchanged.
\textbf{These associations are exploratory and are reported as a candidate regularity only}: with
ten assets a single $p$-value is not decisive, the jump-ratio feature is constant across assets at
the median (one exceedance in 95 returns) and therefore carries no discriminative information at
the pre-registered threshold, and trend persistence flips sign under LOO and is marked
single-asset sensitive.

Two further observations qualify the table. First, the four assets whose confidence intervals
contain zero (USDJPY, ETHUSD, DJI, TM) are all directionally negative and agree in sign with the
supported negative cases, so the negative side is internally consistent even where it is not
individually significant. Second, the two supported positive cases (NVO, BABA) are precisely the
assets with the strongest first-order autocorrelation, which is what the candidate regularity
predicts; the association is therefore not driven by the significance pattern alone but by the
ordering of the effect sizes across the feature range.

\section{Diagnosis and Decisions}
\begin{table}[!t]
\caption{Consolidated answers to the five diagnostic questions.}
\label{tab:q}
\centering
\small
\begin{tabular}{@{}p{0.46\columnwidth}p{0.48\columnwidth}@{}}
\toprule
Question & Answer \\
\midrule
Q1 Robust sample-level routing evidence? &
No. Neither the routing proxy nor the capacity sweep produces a stable state-to-operator mapping. \\
Q2 Genuine asset-level operator heterogeneity? &
Yes. At adequate capacity ETHUSD prefers the nonlinear operator in 36/36 observations, BTCUSD
prefers linear at every capacity, GSPC is near parity. \\
Q3 Explainable by the existing structural features? &
Only as a candidate: realized volatility and $|\mathrm{ACF}_1|$ associate with the interaction sign
and survive LOO, but the sign spans zero across ten assets. \\
Q4 Principal remaining uncertainty? &
Asset dependence and operator form --- not the structural representation, and not routing, which
is now largely excluded. \\
Q5 May a stronger nonlinear operator be reconsidered? &
Only as one candidate, tested against a capacity-matched MLP, and never by extrapolating
``ETHUSD prefers nonlinear'' to a specific operator family. \\
\bottomrule
\end{tabular}
\end{table}

The decisions recorded here are: stop the adaptive-fusion line, since a volatility gate measures
worse than parameter-free addition and the structure it would exploit leaves no switching
headroom; do not build a neural router, because the state-to-operator mapping is not established;
treat operator choice as an asset-level configuration question, which is what the evidence
supports; and keep any structured nonlinear operator conditional on a reproducible context and on
a capacity-matched comparison.

\section{Limitations}
\emph{Operator fidelity:} the proxy operators (a linear map and an MLP over flattened windows)
differ from the framework branches, which combine a DLinear decomposition with an
LSTM--Transformer residual under RevIN. Proxy results must therefore not be read as
``nonlinearity is useless''; within the framework the residual branch still outperforms the linear
base on the S\&P~500 at $T=24$ across both seeds, and the two statements are consistent.
\emph{Sample size:} the asset-level association rests on ten assets and one seed per asset (two
for GSPC), so association intervals are wide and the analysis is exploratory.
\emph{Regime definition:} regimes are volatility quintiles of a causal window measure, and the
calm strata are statistically indistinguishable between operators, so ``calm'' is an acceptance
result rather than a separate effect. \emph{Environment:} runs are CPU-only, so wall-clock times
are not comparable to published GPU numbers and only relative comparisons are interpreted.

\section{Conclusion}
The reproduction reproduces the published ordering at $T=24$ and does so robustly across seeds.
The diagnostics that follow refine but do not overturn it: the marginal value of the nonlinear
branch is asset-dependent rather than regime-universal, the volatility gate does not pay, and
matched-capacity comparisons show that the earlier uniform preference for the linear operator was
partly a capacity artefact. No robust experimental basis for sample-level routing was found,
whereas genuine asset-level operator heterogeneity was. The next question the evidence supports is
therefore not how a router should choose, but in which reproducible contexts a stronger nonlinear
operator is worth its capacity --- to be answered against a capacity-matched baseline. All
artifacts reported here (analysis scripts, batch runners, per-run prediction files and
stratification outputs) are committed to the public repository, so every number in this report
can be regenerated from the committed commands without access to the original authors' data.

\section*{Reproducibility Notes}
All reproduction assets live under \texttt{reproduction/} in the project repository: analysis
scripts under \texttt{reproduction/analysis/}, batch runners under
\texttt{reproduction/batches/}, per-setting results under \texttt{reproduction/results/}. Four
environment facts cost real time and are recorded for reuse: the data pipeline imports
\texttt{patool}, \texttt{huggingface\_hub}, \texttt{sktime} and \texttt{datasets}
unconditionally; \texttt{joblib} must be pinned to \texttt{1.5.3}; CPU runs require
\texttt{--no\_use\_gpu}; and killing DataLoader worker processes (\texttt{spawn\_main})
deadlocks the parent trainer, so a stalled run must be detected by process CPU time rather than
by log output.

\begin{thebibliography}{1}
\bibitem{hsieh2026} C.-Y. Hsieh and Y.-C. Chen, ``DeReFusion: a decomposition-reconstruction
fusion framework for financial time series forecasting,'' \emph{Applied Soft Computing},
vol.~203, p.~116252, 2026.
\bibitem{zeng2023} A. Zeng, M. Chen, L. Zhang, and Q. Xu, ``Are transformers effective for time
series forecasting?'' in \emph{Proc. AAAI}, vol.~37, no.~9, 2023, pp.~11121--11128.
\end{thebibliography}

\end{document}
